% ===========================================================================
\section{Entropy: average surprise}\label{sec:entropy}
\warmth{0}\quad\whisper{Entropy averages the surprise of all possible outcomes.}

\begin{strand}
The surprise $-\log p(x)$ belongs to one outcome. Entropy takes its expectation,
so it depends on the \emph{law} of $X$, not on which value happens to be observed.
\end{strand}

\begin{definition}[entropy in base $b$]\label{def:entropy}
For a discrete random variable $X$ with alphabet $\X$,
\[
  \Ent_b(X)
  =-\sum_{x\in\X}p_X(x)\log_b p_X(x)
  =\E\!\left[-\log_b p_X(X)\right].
\]
We use the convention $0\log_b0=0$. Unless stated otherwise, $b=2$.
\end{definition}

\begin{remark}
Because entropy depends only on $p_X$, the notations
$\Ent(X)$, $\Ent(p_X)$, and $\Ent(p)$ may describe the same number.
\end{remark}

\block{A Bernoulli random variable}

\begin{example}[binary entropy]
Let $\X=\{H,T\}$ with $\PP(X=H)=p$ and $\PP(X=T)=1-p$. Then
\[
  \Ent(X)=-p\log p-(1-p)\log(1-p).
\]
In this special case we sometimes write $\Ent(p)$.
\end{example}

\begin{yourturn}
Use the formula to fill the three anchor values:
\[
  \Ent(0)=\answerin[1cm]{0},\qquad
  \Ent\!\left(\tfrac12\right)=\answerin[1cm]{1},\qquad
  \Ent(1)=\answerin[1cm]{0}\quad\text{bits}.
\]
Sketch the curve for $0\le p\le1$ and mark its highest point.
\workspace[3]
\recall{The values are $0,1,0$. The curve is symmetric about $p=1/2$ and peaks there.}
\end{yourturn}

\block{Joint entropy}

\begin{definition}[joint entropy]\label{def:joint-entropy}
For $X=(X_1,X_2)$ with joint pmf $p_{X_1X_2}$,
\[
  \Ent(X_1,X_2)
  =-\sum_{x_1\in\X_1}\sum_{x_2\in\X_2}
    p_{X_1X_2}(x_1,x_2)\log p_{X_1X_2}(x_1,x_2).
\]
\end{definition}

\begin{proposition}[additivity under independence]\label{prop:independent-entropy}
If $X_1$ and $X_2$ are independent, then
\[
  \Ent(X_1,X_2)=\Ent(X_1)+\Ent(X_2).
\]
\end{proposition}

\begin{proof}
Independence gives
$p_{X_1X_2}(x_1,x_2)=p_{X_1}(x_1)p_{X_2}(x_2)$.
Insert this into Definition~\ref{def:joint-entropy} and split the logarithm:
\begin{align*}
  \Ent(X_1,X_2)
  &=-\sum_{x_1,x_2}p_{X_1}(x_1)p_{X_2}(x_2)
    \bigl(\log p_{X_1}(x_1)+\log p_{X_2}(x_2)\bigr)\\
  &=-\sum_{x_1}p_{X_1}(x_1)\log p_{X_1}(x_1)
      \sum_{x_2}p_{X_2}(x_2)\\
  &\quad-\sum_{x_2}p_{X_2}(x_2)\log p_{X_2}(x_2)
      \sum_{x_1}p_{X_1}(x_1)\\
  &=\Ent(X_1)+\Ent(X_2),
\end{align*}
because each marginal pmf sums to $1$.
\end{proof}

\begin{yourturn}
Recall the key step in the additivity proof:
\[
  \log\!\bigl(p_{X_1}(x_1)p_{X_2}(x_2)\bigr)
  =\answerin[5.2cm]{\log p_{X_1}(x_1)+\log p_{X_2}(x_2)}.
\]
\recall{$\log p_{X_1}(x_1)+\log p_{X_2}(x_2)$.}
\end{yourturn}
